%% Career-Ops LaTeX CV Template
%% Style: ATS-optimized single-column, minimal color accent, single-page friendly.
%% Build: pdflatex {file}.tex (or tectonic {file}.tex)

\documentclass[letterpaper,10pt]{article}

\usepackage[margin=0.55in]{geometry}
\usepackage[T1]{fontenc}
\usepackage[utf8]{inputenc}
\usepackage{lmodern}
\usepackage{microtype}
\usepackage{enumitem}
\usepackage{titlesec}
\usepackage{hyperref}
\usepackage{xcolor}
\usepackage{parskip}

\definecolor{accent}{HTML}{1F3A5F}

\hypersetup{
  colorlinks=true,
  urlcolor=accent,
  linkcolor=accent,
  pdftitle={ Ahnaf An Nafee - Resume },
  pdfauthor={ Ahnaf An Nafee },
  pdfsubject={Resume},
  pdfkeywords={ Site Reliability Engineer, SRE, Kubernetes, OpenShift, Terraform, AWS, Docker, CI/CD, RBAC, observability, infrastructure-as-code, multi-region, high availability, fault tolerance, on-call, incident response, MTTR, Python, Go, Bash, MLOps, distributed systems }
}

% Ensure generated PDF is machine readable / ATS parseable
\input{glyphtounicode}
\pdfgentounicode=1

\pagestyle{empty}
\setlength{\parindent}{0pt}
\raggedbottom
\raggedright

% --- Section formatting ------------------------------------------------------
% Note: section titles render in mixed case (no \MakeUppercase) for maximum
% ATS compatibility. Older parsers normalize headers like "Work Experience"
% and "Education" more reliably than uppercase variants.
\titleformat{\section}
  {\color{accent}\large\bfseries}
  {}{0em}{}[{\color{accent}\vspace{-6pt}\hrulefill\vspace{-4pt}}]
\titlespacing*{\section}{0pt}{8pt}{4pt}

\titleformat{\subsection}
  {\normalsize\bfseries}
  {}{0em}{}
\titlespacing*{\subsection}{0pt}{4pt}{2pt}

% --- Custom commands ---------------------------------------------------------
\newcommand{\role}[4]{%
  \textbf{#1} \hfill \textit{#3}\\[-1pt]
  \textit{#2} \hfill \textit{#4}\par\vspace{2pt}
}

\newcommand{\education}[4]{%
  \textbf{#1} \hfill \textit{#3}\\[-1pt]
  \textit{#2} \hfill \textit{#4}\par\vspace{2pt}
}

\newlist{tightitemize}{itemize}{1}
\setlist[tightitemize]{leftmargin=1.3em, topsep=2pt, itemsep=1pt, parsep=0pt, label=\textbullet}

% =============================================================================
\begin{document}

% --- Header ------------------------------------------------------------------
\begin{center}
  {\LARGE\bfseries\color{accent} Ahnaf An Nafee }\\[3pt]
  {\normalsize Site Reliability Engineer \textbar{} DevOps \& Cloud Infrastructure \textbar{} Applied AI Researcher }\\[4pt]
  \small
  \href{mailto:ahnafnafee@gmail.com}{ahnafnafee@gmail.com} \textbar{}
  540-252-8738 \textbar{}
  Fairfax, VA\\
  \href{https://linkedin.com/in/ahnafnafee}{linkedin.com/in/ahnafnafee} \textbar{}
  \href{https://github.com/ahnafnafee}{github.com/ahnafnafee} \textbar{}
  \href{https://www.ahnafnafee.dev/}{ahnafnafee.dev} \textbar{} \href{https://scholar.google.com/citations?user=u15DO0cAAAAJ}{Google Scholar}
\end{center}

% --- Summary -----------------------------------------------------------------
\section*{Summary}
Site Reliability and DevOps engineer with hands-on production experience running Kubernetes and OpenShift at enterprise scale, automating CI/CD pipelines, and architecting cloud-native infrastructure on AWS. Delivered measurable reliability outcomes: eliminated 100\% of TLS-certificate-related downtime through automated lifecycle management, standardized Kubernetes deployments across 10+ teams, and cut AWS spend by 80\% via right-sizing, autoscaling, and serverless adoption. Comfortable on-call, partnering with security on RBAC and infrastructure hardening, and treating everything-as-code with Terraform and GitHub Actions. PhD candidate in Computer Science (DCXR Lab) bringing applied AI/ML research orientation directly relevant to running reliable training and inference infrastructure for frontier model labs.

% --- Skills ------------------------------------------------------------------
\section*{Technical Skills}
\textbf{Cloud \& DevOps:} AWS, Kubernetes, OpenShift, Docker, Terraform, CI/CD, GitHub Actions, Jenkins, Gradle, Maven, Helm, Observability\\
\textbf{Reliability \& Security:} RBAC, TLS Automation, High Availability, Fault Tolerance, Incident Response, MTTR, Infrastructure-as-Code, System Design\\
\textbf{Languages:} Python, Go, Bash, Java, C++, TypeScript, JavaScript, SQL, GLSL\\
\textbf{AI / ML:} PyTorch, TensorFlow, Generative AI, Deep Learning, LLMs, Diffusion Models, MLOps\\
\textbf{Backend \& Systems:} Distributed Systems, Microservices, REST, gRPC, WebSockets, OAuth\\
\textbf{Methods:} Agile, Scrum, Test-Driven Development, Code Review, Technical Leadership, Cross-Functional Collaboration

% --- Experience --------------------------------------------------------------
\section*{Work Experience}
\role{DevOps Engineer}{Paychex}{Feb 2023 -- Aug 2025}{Rochester, NY}
\begin{tightitemize}
  \item Automated end-to-end TLS certificate lifecycle management on Kubernetes/OpenShift, eliminating manual renewals and reducing certificate-expiration-related production downtime by \textbf{100\%}
  \item Architected and led cross-team standardization of Kubernetes/OpenShift deployment dictionaries across 10+ service teams, eliminating configuration drift and accelerating delivery velocity
  \item Designed and rolled out Role-Based Access Control (RBAC) policies in OpenShift in partnership with security and infrastructure teams, hardening the platform against privilege-escalation risks
  \item Built a Gradle plugin for container image certification that streamlined CI/CD pipelines, reduced build-pipeline congestion, and lowered operational expenses by \textbf{10\%} annually
  \item Championed infrastructure-as-code and observability best practices; partnered with SREs on incident response, deployment reliability, and mean time to recovery (MTTR) improvements
\end{tightitemize}

\role{Chief Technology Officer}{Dynasty 11 Studios}{Jun 2022 -- Feb 2023}{Wayne, PA}
\begin{tightitemize}
  \item Spearheaded migration of monolithic backend to a microservices architecture on AWS, cutting application load and cloud spend by \textbf{80\%} through right-sizing, autoscaling, and serverless adoption
  \item Introduced CI/CD with GitHub Actions and Maven, automating builds, tests, and deployments; reduced manual release effort by \textbf{85\%} and eliminated deploy-day toil
  \item Owned end-to-end technical strategy, architecture, and execution as CTO; led engineering across backend, frontend, and cloud teams shipping a consumer-facing matchmaking platform
  \item Recruited, mentored, and led multiple student engineering pods through full microservice design, implementation, and production rollout
\end{tightitemize}

\role{Software Engineer}{Dynasty 11 Studios}{Sep 2021 -- Jun 2022}{Wayne, PA}
\begin{tightitemize}
  \item Engineered a real-time chat subsystem using STOMP over WebSockets in Java, supporting thousands of concurrent users with low-latency delivery
  \item Integrated 20+ REST endpoints across third-party and OAuth-based identity providers, enabling secure, scalable authentication and data exchange
  \item Optimized client-side state management with Redux, cutting chat-service performance bottlenecks by \textbf{80\%}
  \item Designed data pipelines and APIs powering the matchmaking recommendation engine serving thousands of active users
\end{tightitemize}

\role{Graduate Teaching Assistant}{George Mason University --- DCXR Lab}{Aug 2025 -- Present}{Fairfax, VA}
\begin{tightitemize}
  \item Mentor 60+ graduate and undergraduate students in algorithms, systems, and applied machine learning, driving measurable improvements in lab completion and conceptual mastery
  \item Design coursework, lab exercises, and evaluation rubrics aligned with modern software engineering best practices
\end{tightitemize}

% --- Education ---------------------------------------------------------------
\section*{Education}
\education{Ph.D. in Computer Science}{George Mason University --- DCXR Lab}{}{Fairfax, VA}
\begin{tightitemize}
    \item Research focus: generative AI for 3D content creation, machine-learning-driven graphics pipelines, and human--computer interaction in immersive XR environments
\end{tightitemize}

\education{B.S. in Computer Science}{Drexel University}{}{Philadelphia, PA}
\begin{tightitemize}
    \item Honors: Dean's List; Founder's Scholarship
    \item Relevant coursework: Algorithms, Operating Systems, Distributed Computing, Computer Graphics, Machine Learning, Software Engineering
\end{tightitemize}

% --- Projects ----------------------------------------------------------------
% Section header is "Selected Projects" (not "Research & Selected Projects")
% so ATS classify the items below as Projects. Research focus stays in the
% bullet text itself.
\section*{Selected Projects}
\begin{tightitemize}
\item \textbf{AI-Driven 3D Content Generation} --- Researching generative models (diffusion, neural fields) for automated, controllable 3D asset creation in production-grade graphics pipelines
\item \textbf{ML for Graphics Pipelines} --- Exploring learning-based approaches to UV mapping, stylization, and non-photorealistic rendering (NPR) for real-time applications
\item \textbf{Portfolio \& Publications} --- Additional projects, technical artwork, and publications available at \href{https://www.ahnafnafee.dev/portfolio}{ahnafnafee.dev/portfolio}
\end{tightitemize}

\end{document}
